\section{Affine rays and frozen arithmetic signs}
\label{sec:rays-and-signs}

The divisor coordinates determine exact congruence classes for the two
physical times.  The following statements concern physical solutions:
primality, largest-prime, interval, native-key, and mask conditions may
delete points from the displayed affine rays, but cannot create points
off them.

\subsection{The anchor ray}

\begin{theorem}[Fixed-\(t\) anchor ray]
\label{thm:anchor-ray}
Fix \(m,n,h_0,d_m,d_n\) and a divisor \(t\) supplied by
\cref{thm:anchor-compression}.  If one physical anchor solution
\((j_{R,0},P_{m,0},P_{n,0})\) exists, then every other physical anchor
solution with the same \(t\) is uniquely of the form
\begin{align}
 j_R&=j_{R,0}+abt\,z,
 \label{eq:anchor-time-ray}\\
 P_m&=P_{m,0}+ma\,z,
 \label{eq:anchor-Pm-ray}\\
 P_n&=P_{n,0}+nb\,z
 \label{eq:anchor-Pn-ray}
\end{align}
for an integer \(z\).
\end{theorem}

\begin{proof}
Take two fixed-\(t\) anchor solutions and write differences with a
\(\delta\).  From
\[
 mj_R+h_0=btP_m
\]
one obtains
\[
 m\,\delta j_R=bt\,\delta P_m.
\]
By \cref{rem:primitive-gcd}, \((m,bt)=1\), so \(bt\mid\delta j_R\).
Likewise \(at\mid\delta j_R\).  Since \(a,b,t\) are pairwise coprime,
\[
 \operatorname{lcm}(at,bt)=abt.
\]
Thus \(\delta j_R=abt\,z\).  Substitution gives
\(\delta P_m=ma\,z\) and \(\delta P_n=nb\,z\).  Uniqueness is immediate.
\end{proof}

\begin{corollary}[Anchor census]
\label{cor:anchor-census}
If the allowed anchor-time interval has length \(J_R\), then for fixed
\(t\)
\begin{equation}
 \#\{j_R\text{ in the physical carrier with this }t\}
 \le 1+\frac{J_R}{abt}.
 \label{eq:anchor-census}
\end{equation}
\end{corollary}

\subsection{The erasure ray}

\begin{theorem}[Fixed-\((q,u,v)\) erasure ray]
\label{thm:erasure-ray}
Fix \(m,n,h_0\) and a primitive erasure triple \((q,u,v)\) from
\cref{thm:erasure-divisor-line}.  If one physical erasure solution
\((j_{E,0},Q_{m,0},Q_{n,0})\) exists, then every other physical erasure
solution with the same \((q,u,v)\) is uniquely of the form
\begin{align}
 j_E&=j_{E,0}+quv\,z,
 \label{eq:erasure-time-ray}\\
 Q_m&=Q_{m,0}+mv\,z,
 \label{eq:erasure-Qm-ray}\\
 Q_n&=Q_{n,0}+nu\,z
 \label{eq:erasure-Qn-ray}
\end{align}
for an integer \(z\).
\end{theorem}

\begin{proof}
For two solutions,
\[
 m\,\delta j_E=qu\,\delta Q_m,\qquad
 n\,\delta j_E=qv\,\delta Q_n.
\]
By \cref{rem:primitive-gcd}, \((m,qu)=(n,qv)=1\).  Hence
\(qu\mid\delta j_E\) and \(qv\mid\delta j_E\).  The integers \(q,u,v\)
are pairwise coprime, so their least common multiple is \(quv\).
Writing \(\delta j_E=quv\,z\) gives the two prime increments.
\end{proof}

\begin{corollary}[Erasure census]
\label{cor:erasure-census}
If the allowed erasure-time interval has length \(J_E\), then
\begin{equation}
 \#\{j_E\text{ in the physical carrier with fixed }q,u,v\}
 \le 1+\frac{J_E}{quv}.
 \label{eq:erasure-census}
\end{equation}
\end{corollary}

\subsection{The sign theorem}

\begin{theorem}[Primitive cofactor sign]
\label{thm:primitive-cofactor-sign}
The residual sign in the literal rectangle contribution
\eqref{eq:rectangle-residual-sign} is
\begin{equation}
 \boxed{
 \mu(d_mS_m)\mu(d_nS_n)
 =
 \mu(a)\mu(b)\mu(u)\mu(v).}
 \label{eq:primitive-cofactor-sign}
\end{equation}
Equivalently, the product of the four target M\"obius signs is
\begin{equation}
 \boxed{
 \mu(A_m)\mu(A_n)\mu(B_m)\mu(B_n)
 =
 \mu(a)\mu(b)\mu(u)\mu(v).}
 \label{eq:four-target-sign}
\end{equation}
\end{theorem}

\begin{proof}
Using the squarefree and coprimality conclusions of
\cref{thm:erasure-divisor-line},
\begin{align*}
 \mu(d_mS_m)\mu(d_nS_n)
 &=
 \mu(gaqu)\mu(gbqv)\\
 &=
 \mu(g)^2\mu(q)^2
 \mu(a)\mu(b)\mu(u)\mu(v),
\end{align*}
which proves \eqref{eq:primitive-cofactor-sign}.

For the target signs, each largest prime is coprime to its cofactor and
contributes one minus sign:
\[
 \mu(A_m)\mu(A_n)\mu(B_m)\mu(B_n)
 =
 \mu(R_m)\mu(R_n)\mu(S_m)\mu(S_n).
\]
The anchor factors give
\[
 \mu(bt)\mu(at)=\mu(a)\mu(b),
\]
and the erasure factors give
\[
 \mu(qu)\mu(qv)=\mu(u)\mu(v).
\]
This proves \eqref{eq:four-target-sign}.
\end{proof}

\begin{corollary}[Frozen sign on every prime ray]
\label{cor:frozen-ray-sign}
For fixed rows \(m,n\) and fixed primitive erasure coordinates
\((q,u,v)\), the M\"obius sign of every physical point on
\eqref{eq:erasure-time-ray}--\eqref{eq:erasure-Qn-ray} is the same.
There is no arithmetic sign oscillation in the ray variable \(z\).
\end{corollary}

\begin{remark}[What can still oscillate]
\label{rem:literal-phase-can-oscillate}
The literal amplitude \(\cA\) can vary with the complete native key and
with \(z\).  Its complex phases may cancel.  The theorem says only that
such cancellation cannot be attributed to a varying target M\"obius sign
inside one fixed \((q,u,v)\)-ray.
\end{remark}
